% =========================================================================
% CSR v3 journal draft -- Section I (Introduction) and Section II (Background)
% Direction B (docs/v3/22_V3_PAPER_PLAN.md): v3 CGRA story, journal register.
% Cite keys follow docs/v3/paper/v3_additions.bib; map legacy v2 keys as noted.
% Honesty rules: no result numbers appear in these two sections; every
% architectural claim here is leaf-evidenced (M1-M7) or clearly scoped.
% =========================================================================

\section{Introduction}
\label{sec:introduction}

Compressive sensing (CS) acquires sparse or compressible signals from far
fewer measurements than Nyquist sampling requires
\cite{donoho2006compressed, candes2008introduction}. The price of the reduced
acquisition cost is paid at the receiver: sparse recovery is an iterative,
numerically delicate optimization whose complexity dominates the overall
system energy and latency budget. This asymmetry matters most precisely where
CS is most attractive---edge sensing, implantable and wearable biomedical
interfaces, and embedded communication receivers---where the reconstruction
engine must operate under tight power and real-time constraints rather than
on a workstation.

% Fig. 1 -- Algorithm/application landscape and the runtime operating envelope.
% Requires in preamble:
%   \usepackage{tikz}
%   \usetikzlibrary{arrows.meta,positioning,calc,fit,backgrounds}
% Include from the main text with: % Fig. 1 -- Algorithm/application landscape and the runtime operating envelope.
% Requires in preamble:
%   \usepackage{tikz}
%   \usetikzlibrary{arrows.meta,positioning,calc,fit,backgrounds}
% Include from the main text with: % Fig. 1 -- Algorithm/application landscape and the runtime operating envelope.
% Requires in preamble:
%   \usepackage{tikz}
%   \usetikzlibrary{arrows.meta,positioning,calc,fit,backgrounds}
% Include from the main text with: \input{fig_algorithm_landscape}

\begin{figure}[t]
  \centering
  \resizebox{\columnwidth}{!}{%
  \begin{tikzpicture}[
      font=\scriptsize,
      alg/.style={circle, draw=black, fill=white, minimum size=5.5pt,
                  inner sep=0pt},
      alglabel/.style={inner sep=1pt},
      app/.style={draw=gray!70, fill=gray!8, rounded corners=1pt,
                  inner sep=2pt, text=black!75, align=center},
      appline/.style={gray!70, shorten >=1.5pt},
      traj/.style={-{Latex[length=2mm]}, very thick, red!70!black, dashed},
      >={Latex[length=1.6mm]}
    ]

    % ---- axes ----------------------------------------------------------
    \draw[->] (0,0) -- (8.6,0);
    \node[anchor=north east] at (8.6,-0.12)
      {Per-iteration cost (compute, LS work, state)};
    \draw[->] (0,0) -- (0,5.4);
    \node[rotate=90, anchor=south] at (-0.25,2.7)
      {Empirical recovery robustness under noise};

    % ---- algorithm points ----------------------------------------------
    \node[alg] (mp)     at (1.0,1.0) {};
      \node[alglabel, above left=0pt of mp]{MP};
    \node[alg] (iht)    at (1.5,2.1) {};
      \node[alglabel, above left=0pt of iht]{IHT};
    \node[alg] (gp)     at (2.5,2.5) {};
      \node[alglabel, below right=0pt of gp]{GP};
    \node[alg] (gomp)   at (3.3,3.6) {};
      \node[alglabel, above left=0pt of gomp]{gOMP};
    \node[alg] (omp)    at (4.0,3.1) {};
      \node[alglabel, below right=0pt of omp]{OMP};
    \node[alg] (htp)    at (5.0,3.8) {};
      \node[alglabel, below right=0pt of htp]{HTP};
    \node[alg] (sp)     at (6.3,4.3) {};
      \node[alglabel, above left=0pt of sp]{SP};
    \node[alg] (cosamp) at (7.3,4.6) {};
      \node[alglabel, below right=2pt and 1pt of cosamp]{CoSaMP};

    % ---- application anchors -------------------------------------------
    \node[app, anchor=west] (a_wear) at (2.1,0.75)
      {wearable / implantable telemetry};
    \draw[appline] (a_wear.west) -- (mp);
    \draw[appline] (a_wear.north west) -- (iht);

    \node[app, anchor=west, align=left] (a_mri) at (0.6,4.35)
      {CS-MRI:\\volumes (IHT)\\dynamic (OMP)};
    \draw[appline] (a_mri.south) -- (iht);
    \draw[appline] (a_mri.east) -- (omp);

    \node[app] (a_mmw) at (4.6,1.5)
      {one-bit mmWave\\massive MIMO};
    \draw[appline] (a_mmw.west) -- (gp);

    \node[app] (a_ofdm) at (7.1,2.7)
      {low-SNR OFDM\\channel tracking};
    \draw[appline] (a_ofdm.north) -- (sp);
    \draw[appline] (a_ofdm.north east) -- (cosamp);

    % ---- runtime trajectory of one receiver ----------------------------
    \draw[traj] (cosamp) to[bend right=12] (omp);
    \draw[traj] (omp) to[bend right=10] (iht);
    \node[red!70!black, align=left, anchor=south west] at (2.9,5.05)
      {one deployed receiver at run time:
       SNR $\downarrow$ / battery $\downarrow$ $\Rightarrow$ cheaper program};

    % ---- unified-fabric contour ----------------------------------------
    \begin{scope}[on background layer]
      \node[draw=blue!60, dashed, rounded corners=6pt, thick,
            fit=(mp)(iht)(gp)(omp)(gomp)(htp)(sp)(cosamp),
            inner sep=8pt] (fabric) {};
    \end{scope}
    \node[blue!60!black, align=center, anchor=north] (fablabel)
      at (4.3,-0.55)
      {this work: all eight are compiled context programs on one CGRA
       fabric,\\switched per solve --- no bitstream reconfiguration, no
       resynthesis};
    \draw[blue!60, thick, -{Latex[length=1.8mm]}]
      ($(fablabel.north)+(-1.6,0)$) -- ($(fabric.south)+(-1.9,0)$);

  \end{tikzpicture}%
  }
  \caption{Sparse recovery algorithms occupy different points on the
  robustness--cost trade-off, and applications anchor to different points.
  Positions are ordinal rather than to scale: per-iteration cost follows the
  kernel inventory of Table~\ref{tab:kernels} and the per-iteration cycle
  measurements of Section~\ref{sec:results}; robustness follows published
  empirical comparisons \cite{marques2019review} and the measured
  reconstruction SNR of Section~\ref{sec:results}. Because a single deployed
  system also moves across operating points at run time (dashed trajectory),
  algorithm choice is a runtime decision. CSR maps the entire portfolio onto
  one context-programmed fabric.}
  \label{fig:landscape}
\end{figure}


\begin{figure}[t]
  \centering
  \resizebox{\columnwidth}{!}{%
  \begin{tikzpicture}[
      font=\scriptsize,
      alg/.style={circle, draw=black, fill=white, minimum size=5.5pt,
                  inner sep=0pt},
      alglabel/.style={inner sep=1pt},
      app/.style={draw=gray!70, fill=gray!8, rounded corners=1pt,
                  inner sep=2pt, text=black!75, align=center},
      appline/.style={gray!70, shorten >=1.5pt},
      traj/.style={-{Latex[length=2mm]}, very thick, red!70!black, dashed},
      >={Latex[length=1.6mm]}
    ]

    % ---- axes ----------------------------------------------------------
    \draw[->] (0,0) -- (8.6,0);
    \node[anchor=north east] at (8.6,-0.12)
      {Per-iteration cost (compute, LS work, state)};
    \draw[->] (0,0) -- (0,5.4);
    \node[rotate=90, anchor=south] at (-0.25,2.7)
      {Empirical recovery robustness under noise};

    % ---- algorithm points ----------------------------------------------
    \node[alg] (mp)     at (1.0,1.0) {};
      \node[alglabel, above left=0pt of mp]{MP};
    \node[alg] (iht)    at (1.5,2.1) {};
      \node[alglabel, above left=0pt of iht]{IHT};
    \node[alg] (gp)     at (2.5,2.5) {};
      \node[alglabel, below right=0pt of gp]{GP};
    \node[alg] (gomp)   at (3.3,3.6) {};
      \node[alglabel, above left=0pt of gomp]{gOMP};
    \node[alg] (omp)    at (4.0,3.1) {};
      \node[alglabel, below right=0pt of omp]{OMP};
    \node[alg] (htp)    at (5.0,3.8) {};
      \node[alglabel, below right=0pt of htp]{HTP};
    \node[alg] (sp)     at (6.3,4.3) {};
      \node[alglabel, above left=0pt of sp]{SP};
    \node[alg] (cosamp) at (7.3,4.6) {};
      \node[alglabel, below right=2pt and 1pt of cosamp]{CoSaMP};

    % ---- application anchors -------------------------------------------
    \node[app, anchor=west] (a_wear) at (2.1,0.75)
      {wearable / implantable telemetry};
    \draw[appline] (a_wear.west) -- (mp);
    \draw[appline] (a_wear.north west) -- (iht);

    \node[app, anchor=west, align=left] (a_mri) at (0.6,4.35)
      {CS-MRI:\\volumes (IHT)\\dynamic (OMP)};
    \draw[appline] (a_mri.south) -- (iht);
    \draw[appline] (a_mri.east) -- (omp);

    \node[app] (a_mmw) at (4.6,1.5)
      {one-bit mmWave\\massive MIMO};
    \draw[appline] (a_mmw.west) -- (gp);

    \node[app] (a_ofdm) at (7.1,2.7)
      {low-SNR OFDM\\channel tracking};
    \draw[appline] (a_ofdm.north) -- (sp);
    \draw[appline] (a_ofdm.north east) -- (cosamp);

    % ---- runtime trajectory of one receiver ----------------------------
    \draw[traj] (cosamp) to[bend right=12] (omp);
    \draw[traj] (omp) to[bend right=10] (iht);
    \node[red!70!black, align=left, anchor=south west] at (2.9,5.05)
      {one deployed receiver at run time:
       SNR $\downarrow$ / battery $\downarrow$ $\Rightarrow$ cheaper program};

    % ---- unified-fabric contour ----------------------------------------
    \begin{scope}[on background layer]
      \node[draw=blue!60, dashed, rounded corners=6pt, thick,
            fit=(mp)(iht)(gp)(omp)(gomp)(htp)(sp)(cosamp),
            inner sep=8pt] (fabric) {};
    \end{scope}
    \node[blue!60!black, align=center, anchor=north] (fablabel)
      at (4.3,-0.55)
      {this work: all eight are compiled context programs on one CGRA
       fabric,\\switched per solve --- no bitstream reconfiguration, no
       resynthesis};
    \draw[blue!60, thick, -{Latex[length=1.8mm]}]
      ($(fablabel.north)+(-1.6,0)$) -- ($(fabric.south)+(-1.9,0)$);

  \end{tikzpicture}%
  }
  \caption{Sparse recovery algorithms occupy different points on the
  robustness--cost trade-off, and applications anchor to different points.
  Positions are ordinal rather than to scale: per-iteration cost follows the
  kernel inventory of Table~\ref{tab:kernels} and the per-iteration cycle
  measurements of Section~\ref{sec:results}; robustness follows published
  empirical comparisons \cite{marques2019review} and the measured
  reconstruction SNR of Section~\ref{sec:results}. Because a single deployed
  system also moves across operating points at run time (dashed trajectory),
  algorithm choice is a runtime decision. CSR maps the entire portfolio onto
  one context-programmed fabric.}
  \label{fig:landscape}
\end{figure}


\begin{figure}[t]
  \centering
  \resizebox{\columnwidth}{!}{%
  \begin{tikzpicture}[
      font=\scriptsize,
      alg/.style={circle, draw=black, fill=white, minimum size=5.5pt,
                  inner sep=0pt},
      alglabel/.style={inner sep=1pt},
      app/.style={draw=gray!70, fill=gray!8, rounded corners=1pt,
                  inner sep=2pt, text=black!75, align=center},
      appline/.style={gray!70, shorten >=1.5pt},
      traj/.style={-{Latex[length=2mm]}, very thick, red!70!black, dashed},
      >={Latex[length=1.6mm]}
    ]

    % ---- axes ----------------------------------------------------------
    \draw[->] (0,0) -- (8.6,0);
    \node[anchor=north east] at (8.6,-0.12)
      {Per-iteration cost (compute, LS work, state)};
    \draw[->] (0,0) -- (0,5.4);
    \node[rotate=90, anchor=south] at (-0.25,2.7)
      {Empirical recovery robustness under noise};

    % ---- algorithm points ----------------------------------------------
    \node[alg] (mp)     at (1.0,1.0) {};
      \node[alglabel, above left=0pt of mp]{MP};
    \node[alg] (iht)    at (1.5,2.1) {};
      \node[alglabel, above left=0pt of iht]{IHT};
    \node[alg] (gp)     at (2.5,2.5) {};
      \node[alglabel, below right=0pt of gp]{GP};
    \node[alg] (gomp)   at (3.3,3.6) {};
      \node[alglabel, above left=0pt of gomp]{gOMP};
    \node[alg] (omp)    at (4.0,3.1) {};
      \node[alglabel, below right=0pt of omp]{OMP};
    \node[alg] (htp)    at (5.0,3.8) {};
      \node[alglabel, below right=0pt of htp]{HTP};
    \node[alg] (sp)     at (6.3,4.3) {};
      \node[alglabel, above left=0pt of sp]{SP};
    \node[alg] (cosamp) at (7.3,4.6) {};
      \node[alglabel, below right=2pt and 1pt of cosamp]{CoSaMP};

    % ---- application anchors -------------------------------------------
    \node[app, anchor=west] (a_wear) at (2.1,0.75)
      {wearable / implantable telemetry};
    \draw[appline] (a_wear.west) -- (mp);
    \draw[appline] (a_wear.north west) -- (iht);

    \node[app, anchor=west, align=left] (a_mri) at (0.6,4.35)
      {CS-MRI:\\volumes (IHT)\\dynamic (OMP)};
    \draw[appline] (a_mri.south) -- (iht);
    \draw[appline] (a_mri.east) -- (omp);

    \node[app] (a_mmw) at (4.6,1.5)
      {one-bit mmWave\\massive MIMO};
    \draw[appline] (a_mmw.west) -- (gp);

    \node[app] (a_ofdm) at (7.1,2.7)
      {low-SNR OFDM\\channel tracking};
    \draw[appline] (a_ofdm.north) -- (sp);
    \draw[appline] (a_ofdm.north east) -- (cosamp);

    % ---- runtime trajectory of one receiver ----------------------------
    \draw[traj] (cosamp) to[bend right=12] (omp);
    \draw[traj] (omp) to[bend right=10] (iht);
    \node[red!70!black, align=left, anchor=south west] at (2.9,5.05)
      {one deployed receiver at run time:
       SNR $\downarrow$ / battery $\downarrow$ $\Rightarrow$ cheaper program};

    % ---- unified-fabric contour ----------------------------------------
    \begin{scope}[on background layer]
      \node[draw=blue!60, dashed, rounded corners=6pt, thick,
            fit=(mp)(iht)(gp)(omp)(gomp)(htp)(sp)(cosamp),
            inner sep=8pt] (fabric) {};
    \end{scope}
    \node[blue!60!black, align=center, anchor=north] (fablabel)
      at (4.3,-0.55)
      {this work: all eight are compiled context programs on one CGRA
       fabric,\\switched per solve --- no bitstream reconfiguration, no
       resynthesis};
    \draw[blue!60, thick, -{Latex[length=1.8mm]}]
      ($(fablabel.north)+(-1.6,0)$) -- ($(fabric.south)+(-1.9,0)$);

  \end{tikzpicture}%
  }
  \caption{Sparse recovery algorithms occupy different points on the
  robustness--cost trade-off, and applications anchor to different points.
  Positions are ordinal rather than to scale: per-iteration cost follows the
  kernel inventory of Table~\ref{tab:kernels} and the per-iteration cycle
  measurements of Section~\ref{sec:results}; robustness follows published
  empirical comparisons \cite{marques2019review} and the measured
  reconstruction SNR of Section~\ref{sec:results}. Because a single deployed
  system also moves across operating points at run time (dashed trajectory),
  algorithm choice is a runtime decision. CSR maps the entire portfolio onto
  one context-programmed fabric.}
  \label{fig:landscape}
\end{figure}


No single recovery algorithm serves all of these applications, because the
established methods occupy genuinely different points on a
guarantee--cost--robustness trade-off, and each application domain anchors
to its own point, as illustrated in Fig.~\ref{fig:landscape}. In magnetic resonance imaging, orthogonal matching pursuit
(OMP) \cite{tropp2007omp} is adopted for dynamic reconstruction because its
one-atom-per-iteration refinement delivers high per-frame quality at a
computational cost compatible with gated acquisition \cite{usman2011omp_mri},
while transform-domain CS-MRI favors iterative hard thresholding (IHT)
\cite{blumensath2009iht}, whose constant, least-squares-free iteration cost
is predictable for large image volumes \cite{qu2010iht_mri}. In wireless
receivers the constraints invert: sparse multipath channels must be tracked
under low and rapidly varying SNR, so OFDM channel estimators employ
subspace pursuit (SP) \cite{dai2009sp} and CoSaMP \cite{needell2009cosamp},
whose candidate-expansion-and-pruning structure provides noise-robust
support recovery with RIP-backed guarantees \cite{wu2010sp_ofdm}; when the
number of channel taps is only loosely known, generalized OMP (gOMP)
\cite{wang2012gomp} shortens the iteration count by admitting several atoms
per pass. At millimeter-wave frequencies with one-bit converters, gradient
pursuit (GP) \cite{blumensath2008gp} is preferred precisely because it
replaces repeated least-squares solves with a line search, keeping the
per-iteration cost tractable at massive-MIMO dimensions
\cite{kim2019gp_mmwave}. At the resource-constrained extreme---coarse
telemetry from wearable or implantable nodes---matching pursuit (MP)
\cite{mallat1993mp} and IHT remain relevant simply because they are the
cheapest iterations available, and hard-thresholding pursuit (HTP)
\cite{foucart2011htp} recovers much of the lost accuracy at a modest
least-squares cost when the energy budget allows it.

The decisive observation is that this choice is not fixed even within a
single deployed system (dashed trajectory in Fig.~\ref{fig:landscape}). A CS
receiver crosses operating points at run time:
channel SNR swings between the regime where candidate-pruning methods are
worth their cost and the regime where cheap thresholding suffices; effective
sparsity drifts with the environment; battery state trades reconstruction
quality against energy per solve; and a medical imaging pipeline moves
between real-time preview and offline high-fidelity modes. Algorithm
selection is therefore a \emph{runtime} decision, not a design-time one. A
hardware accelerator that freezes one algorithm freezes the entire operating
envelope; provisioning one dedicated accelerator per algorithm multiplies
silicon and memory cost by the size of the algorithm portfolio; and falling
back to software surrenders the energy advantage that motivated CS
acquisition in the first place. What such systems need is one fabric on
which \emph{any} of the eight algorithms can be invoked per solve---switched
by loading a compiled program image in microseconds, with no
reconfiguration of the FPGA bitstream and no resynthesis---at a silicon
cost bounded by the algorithms' shared kernel structure rather than by
their number.

Existing FPGA and ASIC recovery accelerators do not offer this flexibility.
The dominant line of work hard-wires a single algorithm---most often
OMP---and optimizes its least-squares (LS) stage with dedicated QR or
Cholesky engines \cite{ge2019qr_omp, roy2020ge_omp}. These designs deliver
excellent single-algorithm throughput, but their datapaths, solvers, and
controllers encode one recovery flow. Multi-algorithm designs exist but
support only small families of closely related methods \cite{bai2012omp_amp}.
At the other extreme, software on embedded CPUs or GPUs offers full
algorithmic freedom at an energy cost that defeats the purpose of CS in the
first place.

Between these extremes sits a well-studied middle ground: the coarse-grained
reconfigurable architecture (CGRA), which time-multiplexes a small array of
processing elements (PEs) under compiler-generated, cycle-by-cycle
configuration contexts \cite{liu2019cgra_survey}. Statically
modulo-scheduled CGRAs in the ADRES tradition \cite{mei2003adres} achieve
high utilization on regular loop kernels; elastic fabrics such as STRELA
\cite{strela2024} tolerate variable latency at the cost of per-PE handshake
hardware; hybrid designs such as SNAFU \cite{gobieski2021snafu} and RipTide
\cite{gobieski2022riptide} trade a restricted dataflow discipline for
energy-minimal operation; and open frameworks including CGRA-ME
\cite{ragheb2024cgrame2}, OpenCGRA \cite{tan2020opencgra}, Pillars
\cite{guo2020pillars}, and Morpher \cite{wijerathne2022morpher,
karunaratne2017hycube} provide retargetable compilation, RTL generation, and
simulation. Sparse recovery is, on its face, an ideal CGRA workload: eight
algorithms sharing a compact set of kernels---correlation, matrix--vector
products, selection, residual update, and restricted LS refinement---mapped
onto one fabric by changing programs instead of silicon.

Three gaps, however, separate the published CGRA state of the art from a
deployable multi-algorithm CS engine. First, \emph{commit semantics}:
statically scheduled fabrics stall globally on any boundary backpressure,
while elastic fabrics pay area and energy for handshakes that regular sparse
kernels rarely need; neither exploits the compiler's ability to \emph{prove}
which cycles cannot stall. Second, \emph{trust}: every framework above trusts
its mapper absolutely---an image that references a missing resource, an
illegal route, or an unimplemented operation is discovered, at best, in
simulation and, at worst, silently misexecutes; none certifies configuration
images against the hardware's actual capabilities at load time. Third,
\emph{co-design of the operand stream}: CS recovery is dominated by products
with a sensing matrix that need never be stored, yet no CGRA framework
treats deterministic matrix generation, its numeric scaling, and the reuse
of generated symbols as first-class, compiler-visible resources.

This article presents CSR, a context-programmed CGRA for multi-algorithm
compressive sensing recovery that closes these three gaps. Eight recovery
algorithms---MP, OMP, gOMP, CoSaMP, SP, IHT, HTP, and GP---are compiled into
context programs executed by a dual-cluster, homogeneous $4{\times}4{\times}2$
PE fabric with a single shared program counter; the fabric never decodes an
algorithm identity. The specific contributions are:

\begin{enumerate}
  \item \textbf{A guaranteed-commit elastic--static execution model.} The
  fabric is statically modulo-scheduled, with valid/ready elasticity confined
  to stream and shared-resource boundaries. For every context the compiler
  can prove free of stall dependencies it sets a \emph{guaranteed-commit}
  attribute; the sequencer bypasses the elastic commit join for such
  contexts, and an embedded assertion layer falsifies incorrect proofs. This
  occupies a previously unexplored point between fully static
  \cite{mei2003adres, ragheb2024cgrame2} and fully elastic \cite{strela2024}
  execution.
  \item \textbf{Load-time certified context execution.} Every context image
  is checked \emph{at write time}, field by field, against masks generated
  from the same single-source architecture description that generates the
  RTL parameters and the compiler's resource graph; images are finalized by
  an explicit control-flow-closure scan before they become launchable, and
  operations without an implemented hardware owner fail closed instead of
  silently no-operating. Runtime execution therefore carries no legality
  logic on its critical path.
  \item \textbf{A generated sensing-matrix subsystem with symbol capture.}
  Sensing-matrix entries are produced on demand by a counter-based Threefry
  keyed permutation \cite{salmon2011random}, scaled by a runtime register
  rather than a hard-wired constant, and never stored as a full matrix.
  Symbols observed during selection are captured and promoted into a support
  cache, eliminating the dedicated matrix-fill passes of stored-matrix and
  LFSR-based designs \cite{george2014lfsr}.
  \item \textbf{A certified fixed-point refinement contract.} Restricted
  least-squares refinement is performed by a matrix-free conjugate-gradient
  recurrence with an explicit post-quantization certificate, replacing the
  Gaussian-elimination solvers of prior recovery accelerators
  \cite{roy2020ge_omp}; the fixed-point formats (18-bit data, 27-bit solver,
  62-bit accumulation) were selected by a systematic width sweep rather than
  ad hoc choice, and support-set mutation is transactional with atomic
  commit and rollback.
  \item \textbf{A single-source architecture authority.} One machine-readable
  architecture description generates the RTL parameter headers, the
  modulo-routing resource graph used by the mapper, the certifier masks, and
  a cycle-accurate simulator, and records per-operation timing provenance as
  measured, modeled, or unimplemented---so the compiler cannot schedule
  against hardware that does not exist.
\end{enumerate}

A preliminary version of this architecture appeared in our earlier
conference work \cite{csr_v2}, which mapped the same eight algorithms onto a
24-bit fixed-point PE array with LFSR-generated matrices and a
Gaussian-elimination LS stage. The present design is a ground-up revision:
it introduces the guaranteed-commit context machine, load-time
certification, the Threefry-based generated-matrix subsystem with symbol
capture, the certified matrix-free refinement contract with
systematically selected numeric formats, and a compiler chain built on typed
static single assignment and modulo-routing-graph mapping, and it scales the
reconstruction problem from $(N,M,K)=(256,64,16)$ to $(1024,128,32)$.

The remainder of this article is organized as follows.
Section~\ref{sec:background} reviews compressive sensing, the eight recovery
algorithms and their shared kernels, and the CGRA execution models this work
builds on. Section~\ref{sec:architecture} details the CSR architecture;
Section~\ref{sec:compiler} presents the compiler and certification chain;
Section~\ref{sec:numeric} defines the numeric and refinement contract;
Section~\ref{sec:results} reports implementation results;
Section~\ref{sec:related} positions CSR against prior accelerators and CGRA
frameworks; and Section~\ref{sec:conclusion} concludes.

% =========================================================================

\section{Background}
\label{sec:background}

\subsection{Compressive Sensing}
\label{sec:background_cs}

Let $x \in \mathbb{R}^{N}$ be a $K$-sparse or compressible signal and let
$\Phi \in \mathbb{R}^{M \times N}$, $M < N$, be the sensing matrix. The
measurements are
\begin{equation}
  y = \Phi x + n,
  \label{eq:cs_model}
\end{equation}
where $n$ models measurement noise. Recovering $x$ from the underdetermined
system \eqref{eq:cs_model} exploits sparsity
\cite{donoho2006compressed, candes2008introduction}. Recovery guarantees for
the algorithms considered here are stated in terms of the restricted isometry
property (RIP) of $\Phi$. Random sign (Rademacher) ensembles of the kind
generated on the fly in this work satisfy the RIP \emph{with high
probability} for the problem sizes of interest; as in all
generated-matrix designs, a deployed seed is additionally qualified
empirically rather than claimed to satisfy the RIP deterministically.

\subsection{Sparse Recovery Algorithms and Shared Kernels}
\label{sec:background_algorithms}

This work targets greedy and thresholding-based recovery because their
iterative structure suits hardware acceleration \cite{marques2019review}.
The eight supported algorithms fall into four families with distinct
support-management disciplines:

\begin{itemize}
  \item \emph{Sequential greedy selection.} MP updates one coefficient per
  iteration by projection; OMP adds one support index per iteration and
  re-solves a restricted LS problem; gOMP generalizes the selection to a
  group of indices per iteration, shortening the iteration count.
  \item \emph{Candidate expansion and pruning.} CoSaMP and SP speculatively
  admit $2K$ (respectively $K$) candidates, solve a restricted LS problem on
  the enlarged support (up to $3K$ and $2K$ entries, respectively), and prune
  back to $K$; they differ in how the pruned support is re-estimated.
  \item \emph{Iterative thresholding.} IHT takes a gradient step and applies
  a hard threshold; HTP uses the thresholded support but re-estimates the
  retained coefficients by restricted LS.
  \item \emph{Gradient pursuit.} GP advances along the restricted gradient
  direction with an exact line search, avoiding repeated full LS solves.
\end{itemize}

% Fig. 2 -- Shared iteration template with three concrete instantiations.
% Two-column span (figure*): the six-stage pipeline does not fit one column.
% Requires in preamble:
%   \usepackage{amsmath}
%   \usepackage{amssymb}
%   \usepackage{tikz}
%   \usetikzlibrary{arrows.meta,positioning,calc}
% Include from the main text with: % Fig. 2 -- Shared iteration template with three concrete instantiations.
% Two-column span (figure*): the six-stage pipeline does not fit one column.
% Requires in preamble:
%   \usepackage{amsmath}
%   \usepackage{amssymb}
%   \usepackage{tikz}
%   \usetikzlibrary{arrows.meta,positioning,calc}
% Include from the main text with: % Fig. 2 -- Shared iteration template with three concrete instantiations.
% Two-column span (figure*): the six-stage pipeline does not fit one column.
% Requires in preamble:
%   \usepackage{amsmath}
%   \usepackage{amssymb}
%   \usepackage{tikz}
%   \usetikzlibrary{arrows.meta,positioning,calc}
% Include from the main text with: \input{fig_iteration_template}

\begin{figure*}[t]
  \centering
  \resizebox{\textwidth}{!}{%
  \begin{tikzpicture}[
      font=\scriptsize,
      kern/.style={draw, rounded corners=2pt, fill=blue!6, align=center,
                   minimum height=8mm, minimum width=25mm, inner sep=2pt},
      cell/.style={draw=black!55, rounded corners=1pt, fill=white,
                   align=center, minimum height=7mm, minimum width=25mm,
                   inner sep=2pt, font=\tiny},
      skip/.style={draw=black!35, dashed, rounded corners=1pt, fill=gray!6,
                   align=center, minimum height=7mm, minimum width=25mm,
                   inner sep=2pt, font=\tiny, text=black!50},
      rowlab/.style={anchor=east, align=right, font=\scriptsize\bfseries},
      gen/.style={draw, rounded corners=2pt, fill=orange!12, align=center,
                  inner sep=2.5pt, font=\tiny},
      dat/.style={font=\tiny, text=black!60, above, inner sep=1pt},
      flow/.style={-{Latex[length=1.6mm]}, thick},
      loop/.style={-{Latex[length=1.6mm]}, thick, black!60}
    ]

    % ---- column x positions ---------------------------------------------
    \def\dx{2.95}

    % ---- template row -----------------------------------------------------
    \node[rowlab] at (-1.6,0) {template};
    \node[kern] (t1) at (0,0)        {correlation\\$g=\Phi^{T}r$};
    \node[kern] (t2) at (\dx,0)      {selection\\top-$|{\cdot}|$};
    \node[kern] (t3) at (2*\dx,0)    {support\\mutation};
    \node[kern] (t4) at (3*\dx,0)    {restricted LS\\refinement};
    \node[kern] (t5) at (4*\dx,0)    {residual\\update};
    \node[kern] (t6) at (5*\dx,0)    {stop check /\\certificate};

    \draw[flow] (t1) -- (t2) node[dat, midway]{$g$};
    \draw[flow] (t2) -- (t3) node[dat, midway]{$\Lambda$};
    \draw[flow] (t3) -- (t4) node[dat, midway]{$S$};
    \draw[flow] (t4) -- (t5) node[dat, midway]{$\hat{x}_{S}$};
    \draw[flow] (t5) -- (t6) node[dat, midway]{$r$};
    \draw[loop] (t6.north) -- ++(0,4.5mm) -| (t1.north)
      node[pos=0.25, above, black!60, font=\tiny]{next iteration};

    % ---- generated Phi feed ----------------------------------------------
    \node[gen, anchor=east] (phi) at (-1.1,1.15)
      {$\Phi$ entries generated\\on demand (keyed PRNG)};
    \draw[flow, orange!60!black] (phi.east) -| (t1.70);
    \draw[flow, orange!60!black]
      (phi.east) -- ++(2mm,0) |- ($(t5.north)+(0,2.2mm)$) -- (t5.north);

    % ---- OMP row ----------------------------------------------------------
    \node[rowlab] at (-1.6,-1.35) {OMP};
    \node[cell] at (0,-1.35)      {$g=\Phi^{T}r$};
    \node[cell] at (\dx,-1.35)    {$j=\arg\max|g|$\\(1 atom)};
    \node[cell] at (2*\dx,-1.35)  {$S\leftarrow S\cup\{j\}$};
    \node[cell] at (3*\dx,-1.35)  {$\hat{x}_{S}=\mathrm{LS}(\Phi_{S},y)$\\$|S|\le K$};
    \node[cell] at (4*\dx,-1.35)  {$r=y-\Phi_{S}\hat{x}_{S}$};
    \node[cell] at (5*\dx,-1.35)  {$\|r\|\le\varepsilon$\\or $|S|=K$};

    % ---- CoSaMP row --------------------------------------------------------
    \node[rowlab] at (-1.6,-2.30) {CoSaMP};
    \node[cell] at (0,-2.30)      {$g=\Phi^{T}r$};
    \node[cell] at (\dx,-2.30)    {$\Lambda=\text{top-}2K\,|g|$};
    \node[cell] at (2*\dx,-2.30)  {$T=S\cup\Lambda$\\$|T|\le 3K$};
    \node[cell] at (3*\dx,-2.30)  {$\hat{x}_{T}=\mathrm{LS}(\Phi_{T},y)$;\\prune $S=\text{top-}K$};
    \node[cell] at (4*\dx,-2.30)  {$r=y-\Phi_{S}\hat{x}_{S}$};
    \node[cell] at (5*\dx,-2.30)  {support stable\\or budget};

    % ---- IHT row -----------------------------------------------------------
    \node[rowlab] at (-1.6,-3.25) {IHT};
    \node[cell] at (0,-3.25)      {$g=\Phi^{T}r$};
    \node[cell] at (\dx,-3.25)    {$x'=x+\mu g$;\\keep top-$K$};
    \node[cell] at (2*\dx,-3.25)  {$S=\mathrm{supp}(x')$};
    \node[skip] at (3*\dx,-3.25)  {skipped\\(gradient step only)};
    \node[cell] at (4*\dx,-3.25)  {$r=y-\Phi x'$};
    \node[cell] at (5*\dx,-3.25)  {iteration budget\\or $\|r\|\le\varepsilon$};

    % ---- takeaway ----------------------------------------------------------
    \node[align=left, anchor=north west, font=\tiny, text=black!65,
          text width=157mm] at (-2.9,-3.95)
      {Data types are shared: $g\in\mathbb{R}^{N}$, $r\in\mathbb{R}^{M}$,
       $|S|\le 3K$. Each row is one \emph{compiled context program} executing
       on the same fabric: the stages, their hardware owners, and the data
       widths are identical; only selection width, support discipline, and
       the use of restricted LS differ. gOMP follows the OMP row with $L$
       atoms; SP follows CoSaMP with $K$ candidates and re-estimation; HTP
       follows IHT plus the LS stage; GP and MP follow OMP with the LS stage
       replaced by a line-search or projection step.};

  \end{tikzpicture}%
  }
  \caption{One iteration template (top row) covers all eight recovery
  algorithms, shown instantiated for one representative of each major
  family: OMP (sequential greedy), CoSaMP (candidate expansion and pruning),
  and IHT (iterative thresholding). Vertical alignment makes the claim
  precise: every algorithm sequences the same kernels over the same data
  types, differing only in selection width, support discipline, and whether
  restricted least-squares refinement is used. The sensing matrix is
  generated on demand and feeds both matrix--vector kernels; each row
  corresponds to one compiled context program on the same fabric.}
  \label{fig:iteration}
\end{figure*}


\begin{figure*}[t]
  \centering
  \resizebox{\textwidth}{!}{%
  \begin{tikzpicture}[
      font=\scriptsize,
      kern/.style={draw, rounded corners=2pt, fill=blue!6, align=center,
                   minimum height=8mm, minimum width=25mm, inner sep=2pt},
      cell/.style={draw=black!55, rounded corners=1pt, fill=white,
                   align=center, minimum height=7mm, minimum width=25mm,
                   inner sep=2pt, font=\tiny},
      skip/.style={draw=black!35, dashed, rounded corners=1pt, fill=gray!6,
                   align=center, minimum height=7mm, minimum width=25mm,
                   inner sep=2pt, font=\tiny, text=black!50},
      rowlab/.style={anchor=east, align=right, font=\scriptsize\bfseries},
      gen/.style={draw, rounded corners=2pt, fill=orange!12, align=center,
                  inner sep=2.5pt, font=\tiny},
      dat/.style={font=\tiny, text=black!60, above, inner sep=1pt},
      flow/.style={-{Latex[length=1.6mm]}, thick},
      loop/.style={-{Latex[length=1.6mm]}, thick, black!60}
    ]

    % ---- column x positions ---------------------------------------------
    \def\dx{2.95}

    % ---- template row -----------------------------------------------------
    \node[rowlab] at (-1.6,0) {template};
    \node[kern] (t1) at (0,0)        {correlation\\$g=\Phi^{T}r$};
    \node[kern] (t2) at (\dx,0)      {selection\\top-$|{\cdot}|$};
    \node[kern] (t3) at (2*\dx,0)    {support\\mutation};
    \node[kern] (t4) at (3*\dx,0)    {restricted LS\\refinement};
    \node[kern] (t5) at (4*\dx,0)    {residual\\update};
    \node[kern] (t6) at (5*\dx,0)    {stop check /\\certificate};

    \draw[flow] (t1) -- (t2) node[dat, midway]{$g$};
    \draw[flow] (t2) -- (t3) node[dat, midway]{$\Lambda$};
    \draw[flow] (t3) -- (t4) node[dat, midway]{$S$};
    \draw[flow] (t4) -- (t5) node[dat, midway]{$\hat{x}_{S}$};
    \draw[flow] (t5) -- (t6) node[dat, midway]{$r$};
    \draw[loop] (t6.north) -- ++(0,4.5mm) -| (t1.north)
      node[pos=0.25, above, black!60, font=\tiny]{next iteration};

    % ---- generated Phi feed ----------------------------------------------
    \node[gen, anchor=east] (phi) at (-1.1,1.15)
      {$\Phi$ entries generated\\on demand (keyed PRNG)};
    \draw[flow, orange!60!black] (phi.east) -| (t1.70);
    \draw[flow, orange!60!black]
      (phi.east) -- ++(2mm,0) |- ($(t5.north)+(0,2.2mm)$) -- (t5.north);

    % ---- OMP row ----------------------------------------------------------
    \node[rowlab] at (-1.6,-1.35) {OMP};
    \node[cell] at (0,-1.35)      {$g=\Phi^{T}r$};
    \node[cell] at (\dx,-1.35)    {$j=\arg\max|g|$\\(1 atom)};
    \node[cell] at (2*\dx,-1.35)  {$S\leftarrow S\cup\{j\}$};
    \node[cell] at (3*\dx,-1.35)  {$\hat{x}_{S}=\mathrm{LS}(\Phi_{S},y)$\\$|S|\le K$};
    \node[cell] at (4*\dx,-1.35)  {$r=y-\Phi_{S}\hat{x}_{S}$};
    \node[cell] at (5*\dx,-1.35)  {$\|r\|\le\varepsilon$\\or $|S|=K$};

    % ---- CoSaMP row --------------------------------------------------------
    \node[rowlab] at (-1.6,-2.30) {CoSaMP};
    \node[cell] at (0,-2.30)      {$g=\Phi^{T}r$};
    \node[cell] at (\dx,-2.30)    {$\Lambda=\text{top-}2K\,|g|$};
    \node[cell] at (2*\dx,-2.30)  {$T=S\cup\Lambda$\\$|T|\le 3K$};
    \node[cell] at (3*\dx,-2.30)  {$\hat{x}_{T}=\mathrm{LS}(\Phi_{T},y)$;\\prune $S=\text{top-}K$};
    \node[cell] at (4*\dx,-2.30)  {$r=y-\Phi_{S}\hat{x}_{S}$};
    \node[cell] at (5*\dx,-2.30)  {support stable\\or budget};

    % ---- IHT row -----------------------------------------------------------
    \node[rowlab] at (-1.6,-3.25) {IHT};
    \node[cell] at (0,-3.25)      {$g=\Phi^{T}r$};
    \node[cell] at (\dx,-3.25)    {$x'=x+\mu g$;\\keep top-$K$};
    \node[cell] at (2*\dx,-3.25)  {$S=\mathrm{supp}(x')$};
    \node[skip] at (3*\dx,-3.25)  {skipped\\(gradient step only)};
    \node[cell] at (4*\dx,-3.25)  {$r=y-\Phi x'$};
    \node[cell] at (5*\dx,-3.25)  {iteration budget\\or $\|r\|\le\varepsilon$};

    % ---- takeaway ----------------------------------------------------------
    \node[align=left, anchor=north west, font=\tiny, text=black!65,
          text width=157mm] at (-2.9,-3.95)
      {Data types are shared: $g\in\mathbb{R}^{N}$, $r\in\mathbb{R}^{M}$,
       $|S|\le 3K$. Each row is one \emph{compiled context program} executing
       on the same fabric: the stages, their hardware owners, and the data
       widths are identical; only selection width, support discipline, and
       the use of restricted LS differ. gOMP follows the OMP row with $L$
       atoms; SP follows CoSaMP with $K$ candidates and re-estimation; HTP
       follows IHT plus the LS stage; GP and MP follow OMP with the LS stage
       replaced by a line-search or projection step.};

  \end{tikzpicture}%
  }
  \caption{One iteration template (top row) covers all eight recovery
  algorithms, shown instantiated for one representative of each major
  family: OMP (sequential greedy), CoSaMP (candidate expansion and pruning),
  and IHT (iterative thresholding). Vertical alignment makes the claim
  precise: every algorithm sequences the same kernels over the same data
  types, differing only in selection width, support discipline, and whether
  restricted least-squares refinement is used. The sensing matrix is
  generated on demand and feeds both matrix--vector kernels; each row
  corresponds to one compiled context program on the same fabric.}
  \label{fig:iteration}
\end{figure*}


\begin{figure*}[t]
  \centering
  \resizebox{\textwidth}{!}{%
  \begin{tikzpicture}[
      font=\scriptsize,
      kern/.style={draw, rounded corners=2pt, fill=blue!6, align=center,
                   minimum height=8mm, minimum width=25mm, inner sep=2pt},
      cell/.style={draw=black!55, rounded corners=1pt, fill=white,
                   align=center, minimum height=7mm, minimum width=25mm,
                   inner sep=2pt, font=\tiny},
      skip/.style={draw=black!35, dashed, rounded corners=1pt, fill=gray!6,
                   align=center, minimum height=7mm, minimum width=25mm,
                   inner sep=2pt, font=\tiny, text=black!50},
      rowlab/.style={anchor=east, align=right, font=\scriptsize\bfseries},
      gen/.style={draw, rounded corners=2pt, fill=orange!12, align=center,
                  inner sep=2.5pt, font=\tiny},
      dat/.style={font=\tiny, text=black!60, above, inner sep=1pt},
      flow/.style={-{Latex[length=1.6mm]}, thick},
      loop/.style={-{Latex[length=1.6mm]}, thick, black!60}
    ]

    % ---- column x positions ---------------------------------------------
    \def\dx{2.95}

    % ---- template row -----------------------------------------------------
    \node[rowlab] at (-1.6,0) {template};
    \node[kern] (t1) at (0,0)        {correlation\\$g=\Phi^{T}r$};
    \node[kern] (t2) at (\dx,0)      {selection\\top-$|{\cdot}|$};
    \node[kern] (t3) at (2*\dx,0)    {support\\mutation};
    \node[kern] (t4) at (3*\dx,0)    {restricted LS\\refinement};
    \node[kern] (t5) at (4*\dx,0)    {residual\\update};
    \node[kern] (t6) at (5*\dx,0)    {stop check /\\certificate};

    \draw[flow] (t1) -- (t2) node[dat, midway]{$g$};
    \draw[flow] (t2) -- (t3) node[dat, midway]{$\Lambda$};
    \draw[flow] (t3) -- (t4) node[dat, midway]{$S$};
    \draw[flow] (t4) -- (t5) node[dat, midway]{$\hat{x}_{S}$};
    \draw[flow] (t5) -- (t6) node[dat, midway]{$r$};
    \draw[loop] (t6.north) -- ++(0,4.5mm) -| (t1.north)
      node[pos=0.25, above, black!60, font=\tiny]{next iteration};

    % ---- generated Phi feed ----------------------------------------------
    \node[gen, anchor=east] (phi) at (-1.1,1.15)
      {$\Phi$ entries generated\\on demand (keyed PRNG)};
    \draw[flow, orange!60!black] (phi.east) -| (t1.70);
    \draw[flow, orange!60!black]
      (phi.east) -- ++(2mm,0) |- ($(t5.north)+(0,2.2mm)$) -- (t5.north);

    % ---- OMP row ----------------------------------------------------------
    \node[rowlab] at (-1.6,-1.35) {OMP};
    \node[cell] at (0,-1.35)      {$g=\Phi^{T}r$};
    \node[cell] at (\dx,-1.35)    {$j=\arg\max|g|$\\(1 atom)};
    \node[cell] at (2*\dx,-1.35)  {$S\leftarrow S\cup\{j\}$};
    \node[cell] at (3*\dx,-1.35)  {$\hat{x}_{S}=\mathrm{LS}(\Phi_{S},y)$\\$|S|\le K$};
    \node[cell] at (4*\dx,-1.35)  {$r=y-\Phi_{S}\hat{x}_{S}$};
    \node[cell] at (5*\dx,-1.35)  {$\|r\|\le\varepsilon$\\or $|S|=K$};

    % ---- CoSaMP row --------------------------------------------------------
    \node[rowlab] at (-1.6,-2.30) {CoSaMP};
    \node[cell] at (0,-2.30)      {$g=\Phi^{T}r$};
    \node[cell] at (\dx,-2.30)    {$\Lambda=\text{top-}2K\,|g|$};
    \node[cell] at (2*\dx,-2.30)  {$T=S\cup\Lambda$\\$|T|\le 3K$};
    \node[cell] at (3*\dx,-2.30)  {$\hat{x}_{T}=\mathrm{LS}(\Phi_{T},y)$;\\prune $S=\text{top-}K$};
    \node[cell] at (4*\dx,-2.30)  {$r=y-\Phi_{S}\hat{x}_{S}$};
    \node[cell] at (5*\dx,-2.30)  {support stable\\or budget};

    % ---- IHT row -----------------------------------------------------------
    \node[rowlab] at (-1.6,-3.25) {IHT};
    \node[cell] at (0,-3.25)      {$g=\Phi^{T}r$};
    \node[cell] at (\dx,-3.25)    {$x'=x+\mu g$;\\keep top-$K$};
    \node[cell] at (2*\dx,-3.25)  {$S=\mathrm{supp}(x')$};
    \node[skip] at (3*\dx,-3.25)  {skipped\\(gradient step only)};
    \node[cell] at (4*\dx,-3.25)  {$r=y-\Phi x'$};
    \node[cell] at (5*\dx,-3.25)  {iteration budget\\or $\|r\|\le\varepsilon$};

    % ---- takeaway ----------------------------------------------------------
    \node[align=left, anchor=north west, font=\tiny, text=black!65,
          text width=157mm] at (-2.9,-3.95)
      {Data types are shared: $g\in\mathbb{R}^{N}$, $r\in\mathbb{R}^{M}$,
       $|S|\le 3K$. Each row is one \emph{compiled context program} executing
       on the same fabric: the stages, their hardware owners, and the data
       widths are identical; only selection width, support discipline, and
       the use of restricted LS differ. gOMP follows the OMP row with $L$
       atoms; SP follows CoSaMP with $K$ candidates and re-estimation; HTP
       follows IHT plus the LS stage; GP and MP follow OMP with the LS stage
       replaced by a line-search or projection step.};

  \end{tikzpicture}%
  }
  \caption{One iteration template (top row) covers all eight recovery
  algorithms, shown instantiated for one representative of each major
  family: OMP (sequential greedy), CoSaMP (candidate expansion and pruning),
  and IHT (iterative thresholding). Vertical alignment makes the claim
  precise: every algorithm sequences the same kernels over the same data
  types, differing only in selection width, support discipline, and whether
  restricted least-squares refinement is used. The sensing matrix is
  generated on demand and feeds both matrix--vector kernels; each row
  corresponds to one compiled context program on the same fabric.}
  \label{fig:iteration}
\end{figure*}


Despite these differences, the families draw on one compact kernel set,
summarized in Table~\ref{tab:kernels} and visualized as a single iteration
template in Fig.~\ref{fig:iteration}: forward projection $\Phi x$,
correlation $\Phi^{T} r$, residual update, support selection (a top-$|{\cdot}|$
reduction with deterministic tie-breaking), support-set mutation
(append, union, prune, rollback), restricted LS refinement, and---specific
to generated-matrix architectures---on-demand production of the $\Phi$
entries themselves. Two structural observations drive the architecture in
this article. First, every kernel is either a streaming
multiply--accumulate over $\Phi$ rows or columns, or a compact reduction and
bookkeeping operation over at most $3K$ scalars; none requires storing
$\Phi$, provided its entries can be regenerated deterministically. Second,
as Fig.~\ref{fig:iteration} makes explicit, the kernels compose into
algorithms purely through \emph{control sequencing}---which kernel runs
next, at what selection width, on which support view, under which
convergence predicate---which is exactly the part of the computation that a
context-programmed machine can replay from a compiled image instead of
encoding in fixed-function control logic.

% Table~\ref{tab:kernels}: upgrade of the conference Table I. Columns:
% forward, correlation, residual, selection, support mutation, restricted LS,
% pruning, generated-Phi. Rows grouped by the four families above.
% (Table body to be filled from docs/v3/02_ALGORITHM_CONTRACT.md.)

The candidate-expansion methods bound the live state: SP solves restricted
problems of size at most $2K$ and CoSaMP of size at most $3K$ before
pruning. These bounds, not $N$, size the on-fabric support workspace, which
is what allows a fixed hardware configuration to scale to
$(N,M,K)=(1024,128,32)$.

\subsection{CGRA Execution Models}
\label{sec:background_cgra}

A CGRA exposes an array of word-level PEs whose operations and
interconnections are set by configuration contexts
\cite{liu2019cgra_survey}. Two axes of the design space matter for this
work.

\emph{Spatial versus spatio-temporal mapping.} Spatial-only fabrics pin each
dataflow-graph node to one PE for the duration of a kernel, which suits
small kernels and low power but caps the graph size at the array size.
Spatio-temporal fabrics reconfigure every cycle from a context memory,
time-multiplexing large graphs onto a small array under modulo scheduling
with an initiation interval (II) \cite{mei2003adres}; the compiler maps
operations onto a time-unrolled modulo routing resource graph (MRRG)
\cite{ragheb2024cgrame2, wijerathne2022morpher}. CS recovery graphs are far
larger than any practical array, so CSR is spatio-temporal.

\emph{Static versus elastic timing.} Fully static fabrics assume every
latency is compile-time known; any unpredictable boundary event (a memory
bank conflict, a variable-latency divider, DMA interference) must stall the
whole array atomically. Fully elastic fabrics insert valid/ready handshakes
throughout the interconnect \cite{strela2024}, tolerating unknown latencies
at a per-PE hardware cost that regular kernels do not repay; hybrid
disciplines such as ordered dataflow in SNAFU \cite{gobieski2021snafu}
restrict the dynamism to what irregular workloads actually need. CS
recovery kernels are regular, but their \emph{boundaries}---streamed
scratchpad access, a shared divider, generated-matrix delivery---are not.
CSR therefore keeps the fabric statically scheduled, terminates all
elasticity in stream and shared-resource interfaces, and lets the compiler
discharge the residual cost: contexts proven independent of any elastic
boundary commit unconditionally (Section~\ref{sec:architecture}).

One further property separates this design from the frameworks above. In
CGRA-ME, OpenCGRA, Pillars, and Morpher, the mapper's output is trusted: no
runtime mechanism checks a configuration image against the fabric's actual
capabilities. For an accelerator that hosts eight algorithm programs,
accepts field updates of those programs, and shares resources among them,
CSR instead treats the architecture description as a single generated
authority and enforces it twice---once at compile time in the mapper, and
once at image-load time in hardware---so that the executing array never
carries legality checks on its cycle-critical path, yet no uncertified
context can commit architectural state.
